\section{Reporting Patterns and Sensitivity Census}\label{app:r34}
This historical census recomputes six synthesis findings from the frozen record judgments and map, using the extended 20-family inventory that includes AHB's Type III evasion records. The current privacy/attribution cohort and its 19-family projections are in S41. No source cell, pair judgment or experimental observation is changed. The executable census \path{tools/reporting_patterns.py} saves every numerator/denominator member and input hash in \path{data/reporting-patterns.json}. Family availability uses the same B, P, ?, U, A, inapplicable precedence as the main matrix. Fractions describe reports in the inspected source frame. The sensitivity operations change conventions or inventory membership, not source facts or coding reliability.

\subsection{Inference, utility, fitting and matched adaptation}
\begin{table}[t]\centering\footnotesize
\caption{Reporting asymmetry across the 20 mitigation families. Each denominator contains families to which that item applies; B is explicit, P partial. The common-population row removes denominator differences.}\label{tab:r34-reporting}
\begin{tabularx}{\columnwidth}{@{}Yrr@{}}\toprule
Reported binding & B / applicable & (B+P) / applicable\\\midrule
Inference outcome (R7) & 16/17 & 17/17\\
Cost (R9) & 11/19 & 15/19\\
Utility (R8) & 4/17 & 6/17\\
Protected fitting (R4) & 4/11 & 9/11\\
Matched adaptation (R5) & 2/11 & 2/11\\
\midrule
\multicolumn{3}{@{}l}{Same 14 families, inference / cost / utility:}\\
Explicit B & \multicolumn{2}{l}{13/14\quad 6/14\quad 3/14}\\
B or P & \multicolumn{2}{l}{14/14\quad 10/14\quad 5/14}\\

\bottomrule\end{tabularx}
\par\smallskip\raggedright Expanding R4/R5 applicability to all 20 gives 4/20 and 2/20 (B), or 9/20 and 2/20 (B+P). Source judgments stay fixed. Removal checks and member identities accompany the census (main Appendix F; S38).
\end{table}

The common population contains 14 families: KVGov, SparLeak, SafeKV, Dangol, Early Bird, Shape and Substance, WRA, NetEcho, the cloud benchmark, SpliceLeak, Time Will Tell, Whisper Leak, CachePrune and AHB. On these same families, explicit inference/cost/utility counts are 13/6/3, increasing to 14/10/5 when P is accepted. This matched denominator supports the descriptive cost--utility asymmetry without dividing percentages from differently eligible populations.

Protected fitting has five partial bindings, so its count changes from 4/11 to 9/11 under P acceptance. Matched adaptation has two B, two ?, seven U, and no P. Calling fitting itself generally absent would erase relevant partial evidence. Even accepting the two unresolved family R5 bindings optimistically yields only 4/11; this is a bound on the current uncertainty, not an assignment of B. The broad applicability check temporarily treats all nine inapplicable R4/R5 family entries as U: counts become 4/20 and 2/20, or 9/20 and 2/20 with P. This deliberately stricter denominator is not a revised requirement on non-learning attacks or formal guarantees.
\begin{table*}[t]\centering\footnotesize
\caption{Sensitivity of reporting patterns. Fractions give explicit / applicable family counts; bracketed numerators also accept P. Removing final-search additions excludes six records from four families; contextual controls are excluded in every row. IC separate means no single record carries both explicit inference and cost.}\label{tab:r34-removal}
\begin{tabularx}{\textwidth}{@{}Yccccccl@{}}\toprule
Inventory & R7 & R9 & R8 & R4 & R5 & Agent R8 & IC separate\\\midrule
All mitigation & 16/17 [17] & 11/19 [15] & 4/17 [6] & 4/11 [9] & 2/11 [2] & 0/4 [2] & 5/8\\
Without CachePrune & 15/16 [16] & 10/18 [14] & 3/16 [5] & 4/11 [9] & 2/11 [2] & 0/4 [2] & 5/7\\
Without final-search records & 13/13 [13] & 9/15 [12] & 2/14 [4] & 4/9 [8] & 2/9 [2] & 0/4 [2] & 4/6\\
\bottomrule\end{tabularx}\end{table*}


\subsection{Agent utility and polarity}
\begin{table}[ht]\centering\footnotesize
\caption{Agent privacy/attribution subset: P5 generation and Type I or II. Every applicable family is listed.}\label{tab:r34-agent}
\begin{tabularx}{\linewidth}{@{}lclY@{}}\toprule
Family & R8 & Objective & Located utility evidence\\\midrule
UI Traces~\cite{uitraces2026} & U & Attribution & Offline delay; task outcome unbound.\\
WRA~\cite{jeong2026wra} & P & Privacy & Separately judged usefulness; no same-condition task pair.\\
Dangol~\cite{dangol2026workflow} & U & Privacy & Capture/loopback outcome; workflow utility unbound.\\
Cloud benchmark~\cite{pouryousef2026traffic} & P & Privacy/attribution & Semantic preservation asserted; task score unbound.\\
\bottomrule\end{tabularx}\end{table}

The subset rule is P5 generation and Type I or II. Its zero explicit task-utility result becomes two of four under P acceptance, so the zero is specific to explicit binding. WRA's separately rated usefulness and the cloud benchmark's semantic-preservation statement remain partial evidence. AHB's task accuracy is explicitly bound but belongs to Type III detector evasion. TDXRay, CachePrune and SparLeak supply component utility in different settings. Removing CachePrune, the six final-search records, or both leaves the four-member agent subset unchanged. Contextual records never enter this denominator.

\subsection{Coverage of the empirical map}
\begin{table}[t]\centering\footnotesize
\caption{Core-map coverage by source-layer block. Numerators are families with an evaluated mitigation record; denominators are the empirical core map, not all defense research.}\label{tab:r34-coverage}
\begin{tabularx}{\columnwidth}{@{}Yrr@{}}\toprule
Generating-layer block & Mitigation & Core\\\midrule
P1 only & 1 & 8\\
P1/P2 & 4 & 8\\
P2 only & 10 & 22\\
Orchestration/workflow & 4 & 15\\
Other cross-layer & 1 & 5\\
\bottomrule\end{tabularx}
\par\smallskip\raggedright Across overlapping purpose labels, Type II has 2/13 mitigation families (UI Traces and the cloud benchmark). C5 and the scoped evaluation rule define this coverage frame.
\end{table}

The five blocks partition all 58 core families. The numerator requires a scoped evaluated mitigation record, and the denominator requires C1--C5. Thus hardware coverage of 1/8 does not mean only one hardware defense exists; adjacent proof-only and defense-only studies are outside this empirical frame. Holding the 58-family map fixed but excluding final-search mitigation records changes the five block numerators from 1/4/10/4/1 to 1/1/9/4/1; the service-only block remains the largest and the Type II numerator remains two. This removal checks dependence on added mitigation evidence, not completeness of the map. The Type II count is an overlapping purpose-label incidence, not a sixth block. The machine-readable census gives every family identity and its generating/observing layer set.

\subsection{Inference/cost location and measurement mode}
\begin{table}[ht]\centering\footnotesize
\caption{Sensitivity to accepting partial inference and cost evidence. A family is counted only if it has both requested items under the specified rule.}\label{tab:r34-cost-relaxation}
\begin{tabular}{@{}lrrrr@{}}\toprule
Inventory & Both B & Separate & Both B/P & Separate\\\midrule
All mitigation & 8 & 5 & 12 & 3\\
Without CachePrune & 7 & 5 & 11 & 3\\
Without final-search records & 6 & 4 & 9 & 3\\
\bottomrule\end{tabular}\end{table}

Under strict B, the five separated families are Early Bird, SpliceLeak, Time Will Tell, KVGov and OTRO. The result is five of seven after removing CachePrune and four of six after removing the final-search additions. If P also counts, only SpliceLeak, Time Will Tell and KVGov remain separated among 12 eligible families (three of 11 without CachePrune, three of nine without the final-search records). This is why the main finding names explicit bindings and does not claim separation for most partially specified evidence. Full record identities are retained in the census and the main-PDF crosswalk.

The 39-record mode distribution is Ex20/Tr6/Si5/Co6/Un2. All six contextual records are Ex, making the mitigation-only distribution Ex14/Tr6/Si5/Co6/Un2 (33 total). Ex denotes executed collection/prototype, not necessarily end-to-end production execution. The network receiver subset, defined by P3 in the receiving-layer set, has 12 records: WRA; two NetEcho; three Whisper Leak; SPEC; two Dangol; two cloud benchmark; and remote timing. Its six Ex, four Tr and two Un records do not support a claim that network evidence is mainly replay. The final-search exclusion leaves Ex11/Tr6/Si3/Co5/Un2 (27 mitigation records). Six Co records include five cost-bearing components and KVGov's unprotected hardware calibration; none automatically prices an entire protected deployment.

\subsection{Holdout evidence objects}
The archived matrix has 13 task/input, three model, two topology and one network/load H entries; hardware and time have no H entry in that extraction. Four task H families belong to the mitigation audit. The three model H entries (LLM-Fin, ArchSnoop, LeakyLMs) evaluate attack/identification generalization outside the mitigation families. The two topology H entries (Dangol and the cloud benchmark) likewise evaluate workflow recognition. The source reconciliation in S34 gives their exact locators and objects. No frozen source-to-target defense calibration is established by those five entries. This scoped conclusion does not recode the other 14 H entries as uniformly attack-only or treat unresolved hardware/time cells as evidence that no experiment exists. Accepting variation V as coverage still requires identifying what was held fixed or locally refitted. Time Will Tell and Whisper Leak remain useful examples of defense-condition variation.
